\documentclass[10pt,twocolumn]{article}

\usepackage[utf8]{inputenc}
\usepackage[T1]{fontenc}
\usepackage{lmodern}
\usepackage[margin=0.75in]{geometry}
\usepackage{amsmath,amssymb}
\usepackage{booktabs}
\usepackage{array}
\usepackage{multicol}
\usepackage{microtype}
\usepackage[hidelinks]{hyperref}

\setlength{\columnsep}{0.28in}

\title{\textbf{Feature Ablation Study for Deep Learning\\
Statistical Arbitrage on IDX Equities}}
\author{trading-BEI project}
\date{July 2026}

\begin{document}
\maketitle

\begin{abstract}
\noindent
We describe the experimental setup for determining the optimal number
and composition of input features for a deep learning statistical
arbitrage (DLSA) strategy on Indonesia Stock Exchange (IDX) equities.
Unlike the original DLSA pipeline, no factor model is used: normalized
raw features are fed directly into a two-stage cross-sectional
Transformer --- temporal self-attention over each stock's 60-day window,
then self-attention \emph{across} all stocks on a given day --- that
learns representation and trading signal end-to-end. The model is trained
with a DLSA-style economic objective: maximize the annualized Sharpe
ratio of a long-only softmax portfolio, \emph{net} of transaction costs,
over blocks of consecutive trading days so that turnover is genuinely
penalized. Because the deployed strategy is \emph{long-only}, it carries
market beta and is therefore benchmarked against a buy-and-hold IHSG
index. We use a nested, group-wise ablation across five economically
motivated feature groups, run over eight random seeds per configuration,
and evaluated on strictly time-ordered walk-forward out-of-sample data.
The optimal feature set is identified as the point at which the
walk-forward Sharpe ratio saturates before additional features degrade
performance through overfitting.
\end{abstract}

\section{Introduction}
The \texttt{trading-BEI} project adapts Deep Learning Statistical
Arbitrage (Guijarro-Ordonez, Pelger \& Zanotti, 2021) to IDX equities.
The key deviation from upstream DLSA is the removal of the factor model
and residual stages: raw daily stock-summary data is normalized and
passed directly to a Transformer trading policy.

A central open question is how many input features the model should
receive. Because there is no factor model to pre-extract structure, the
Transformer must learn cross-asset co-movement itself. Adding features
increases model capacity but, on a limited history of roughly $1{,}250$
trading days, also increases the risk of overfitting. This report
specifies a controlled ablation study to locate the optimal feature
count and composition.

\section{Objective}
The project has two nested objectives. The \emph{research} objective is
to identify how many and which input features maximize risk-adjusted
out-of-sample performance under realistic IDX trading frictions --- the
ablation reported here. The \emph{model} objective, optimized during
training, is economic rather than statistical: rather than regressing
next-day returns, the network directly maximizes the annualized Sharpe
ratio of the long-only portfolio it induces, net of transaction costs.
Formally, for a block of consecutive rebalance periods the training loss
is $-\mathrm{Sharpe}$ of the net excess-return series (Section~7). This
aligns the training signal with what is actually traded and evaluated,
and --- because costs are charged inside the block --- makes the model
pay for churn. The headline decision criterion for the ablation is the
walk-forward net Sharpe of the traded top-$N$ rule; cumulative return
versus a buy-and-hold IHSG benchmark, maximum drawdown, and turnover are
reported alongside it.

\section{Data}
The raw input is the IDX daily \emph{Ringkasan Saham} (stock summary),
one record per (date, ticker), spanning roughly six years
(2020-07 to 2026-07). After cleaning, the panel retains OHLC prices,
previous close, volume, value, trade frequency, foreign buy/sell, listed
and tradeable shares, and end-of-day bid/offer with sizes. A single
processed \texttt{panel.parquet} feeds all experiments.

\paragraph{Universe.} A causal liquidity screen is applied: a name is
tradable on day $t$ only if its trailing 20-day \emph{median} daily
traded value is at least $1$ billion IDR (\texttt{market.apply\_universe}).
This keeps penny stocks out of both the tradable set and the per-day
z-score statistics.

\paragraph{Corporate actions.} IDX's raw \texttt{close} is not
split-adjusted, but the reported \texttt{Previous} (prev\_close) is
adjusted on the ex-date, so daily returns use $\ln(C_t/\text{prev}_t)$
(falling back to the lagged close), turning a $1{:}10$ split into a
$\approx 0\%$ return rather than a spurious $-90\%$. An adjusted
log-price path drives all multi-day quantities so momentum and the
forward target are split-consistent.

\paragraph{Data-quality guards.} The raw \texttt{OpenPrice} is frequently
zero; open-based and range features mask non-positive prices to
\texttt{NaN} (later neutralized to $0$ after standardization) rather than
dropping rows. Rows are dropped only when the core feature
(\texttt{log\_return}) is missing, so the sample set is invariant to
which optional feature groups are active.

\section{Strategy Setting}
The strategy is \emph{long-only}: positions lie in $[0,1]$ rather than
the dollar-neutral $[-1,1]$ of classical statistical arbitrage. The
model head therefore uses a sigmoid (or clamped \texttt{tanh}), and
weights are normalized to a fully-invested (or cash-permitting) book.

A direct consequence is exposure to market beta: unlike a
market-neutral long/short book, a long-only strategy moves with the
index. The appropriate benchmark is therefore a buy-and-hold IHSG index,
not an absolute Sharpe threshold. We recommend training with the full
cross-sectional signal (allowing the model to learn both out- and
under-performers) and applying the long-only constraint only at the
portfolio-construction stage.

\section{Model Architecture}
The flagship model is a \emph{cross-sectional} Transformer
(\texttt{models/cross\_sectional.py}) that processes one trading day at a
time. Its input is the tensor of all $N$ eligible stocks on that day,
each with a $T{=}60$-day lookback window of $F$ standardized features,
shape $(N, T, F)$; the output is one score per stock, $(N,)$, which is
ranked to select the long book. The forward pass has three stages:

\begin{enumerate}
\item \textbf{Input projection + temporal encoding.} A linear map
$\mathbb{R}^{F}\!\to\!\mathbb{R}^{d}$ plus a learned positional embedding
over the $T$ time steps. A stack of \texttt{temporal\_layers} standard
Transformer encoder layers applies self-attention \emph{over the 60 days}
of each stock independently, yielding a $(N, T, d)$ tensor.
\item \textbf{Pooling.} The last time step is taken (\texttt{pooling:
last}), collapsing each stock to a single $d$-dimensional vector, $(N,
d)$.
\item \textbf{Cross-sectional encoding + head.} The $N$ stock vectors are
treated as an unordered set (no positional encoding) and passed through
\texttt{cross\_layers} Transformer encoder layers of self-attention
\emph{across stocks}, so each stock's representation is informed by the
whole market that day. A \texttt{LayerNorm} + linear head emits one
scalar score per stock.
\end{enumerate}

The cross-sectional stage is the key departure from the per-stock
baseline Transformer (\texttt{models/transformer.py}), which scores each
stock in isolation; here scoring is a \emph{relative} comparison across
the day's cross-section, and there is no factor model --- attention runs
directly on normalized raw features. Table~\ref{tab:arch} lists the
architecture hyperparameters, held fixed across all ablation experiments.

\begin{table}[t]
\centering
\small
\caption{Architecture and optimization settings (locked across A--E).}
\label{tab:arch}
\begin{tabular}{@{}ll@{}}
\toprule
Setting & Value \\
\midrule
Model width $d_{\text{model}}$ & 64 \\
Attention heads & 4 \\
Temporal / cross layers & 2 / 2 \\
Feed-forward dim & 128 \\
Dropout & 0.1 \\
Lookback $T$ / horizon & 60 / 1 \\
Pooling & last step \\
\midrule
Optimizer & Adam \\
Learning rate & $3\!\times\!10^{-4}$ \\
Weight decay & $1\!\times\!10^{-5}$ \\
Epochs (max) & 40 \\
Early-stop patience & 8 \\
Batch granularity & 1 day / step \\
\bottomrule
\end{tabular}
\end{table}

\section{Training Objective}
The model is trained with a DLSA-style economic objective
(\texttt{train\_dlsa}) rather than regression. Each gradient step draws a
block of $K{=}32$ \emph{consecutive} trading days. On each day the model
scores every stock; a softmax over the cross-section (temperature
$\tau{=}0.1$, concentrating weight like the traded top-$N$ book) produces
long-only weights $w \in [0,1]^{N}$ that sum to one. An extra always-zero
``cash'' logit (\texttt{allow\_cash}) lets the portfolio retreat to the
risk-free rate. Day-over-day weight changes within the block, aligned by
ticker, are charged asymmetric buy/sell costs, and the loss is the
negative annualized Sharpe of the block's \emph{net} excess returns:
\[
\mathcal{L} = -\,\frac{\overline{r}^{\,\text{net}} - r_f}
{\mathrm{std}(r^{\text{net}})}\,\sqrt{A},
\qquad
r^{\text{net}}_t = w_t^{\!\top} \rho_t - c_t,
\]
where $\rho_t$ are simple returns, $c_t$ the ticker-aligned turnover
cost, and $A$ the number of rebalance periods per year. Model selection
and early stopping use the net Sharpe of the \emph{actually traded} rule
--- top-$N$ equal weight after costs --- on the validation days, so the
model selected is the one truly traded, not the softmax portfolio.

\section{Feature Groups}
Ablation is performed over \emph{groups} of economically related
features rather than individual features, to limit the number of
experiments and reduce noise. Each feature is standardized
cross-sectionally within each trading day, following the DLSA
convention. Table~\ref{tab:groups} lists the five groups.

\begin{table}[t]
\centering
\small
\caption{Feature groups used in the ablation.}
\label{tab:groups}
\begin{tabular}{@{}clc@{}}
\toprule
Grp & Features & \# \\
\midrule
G1 & log\_return, mom\_5/10/20 & 4 \\
G2 & hl\_range, roll\_vol\_20, range\_ratio & 3 \\
G3 & log\_volume, log\_value, turnover, amihud & 4 \\
G4 & foreign\_flow\_ratio, foreign\_roll\_5 & 2 \\
G5 & bid\_ask\_spread, book\_imbalance & 2 \\
\bottomrule
\end{tabular}
\end{table}

Representative feature definitions, all causal (no look-ahead):
\begin{align}
r_t &= \ln(C_t / C_{t-1}) \\
\text{hl\_range}_t &= (H_t - L_t)/C_t \\
\text{turnover}_t &= V_t / \text{shares}_t \\
\text{amihud}_t &= |r_t| / \text{value}_t \\
\text{ff\_ratio}_t &= (FB_t - FS_t)/(V_t + 1)
\end{align}
where $C,H,L$ are close/high/low, $V$ volume, $FB/FS$ foreign buy/sell.

\section{Ablation Design}
Experiments are \emph{nested}: starting from a minimal set and adding one
group at a time (Table~\ref{tab:design}). The optimal feature count is
the stage at which the walk-forward Sharpe stops improving.

\begin{table}[t]
\centering
\small
\caption{Nested experiment design.}
\label{tab:design}
\begin{tabular}{@{}clc@{}}
\toprule
Exp & Active groups & $\approx$\# feat. \\
\midrule
A (baseline) & G1 & 4 \\
B & G1+G2 & 7 \\
C & G1+G2+G3 & 11 \\
D & G1+G2+G3+G4 & 13 \\
E (full) & G1+G2+G3+G4+G5 & 15 \\
\bottomrule
\end{tabular}
\end{table}

If the Sharpe rises through A$\to$B$\to$C then plateaus or declines at
D/E, the saturation point identifies the optimal feature set. An
optional leave-one-group-out pass on the best configuration confirms
each group's marginal contribution.

\section{Experimental Controls}
For the comparison to be valid, \emph{only the active feature list} may
differ between experiments. Every other knob is locked in a shared base
config (\texttt{configs/ablation/\_base.yaml}); experiments A--E inherit
it and override only \texttt{experiment\_name} and
\texttt{feature\_groups}. Held identical are: the architecture and
optimization settings of Table~\ref{tab:arch}; the lookback $T{=}60$ and
horizon; the walk-forward split dates; the long-only, top-$N{=}10$
portfolio construction; the transaction-cost and turnover model; and the
seed grid. Without these controls, differences in Sharpe cannot be
attributed to features.

\paragraph{Walk-forward splits.} Evaluation uses a rolling retrain. The
out-of-sample period begins 2024-07-01 and is covered by four folds of
six months each; every fold trains up to a six-month validation window
placed immediately before its test period, trades the next six months
out-of-sample, then rolls forward. Model weights are re-initialized with
an identical seed per fold so differences come from data, not
initialization, and the four folds' test scores are stitched into one
two-year out-of-sample series before backtesting.

\paragraph{Execution timing.} Features on day $t$ use day-$t$ closing
data, so orders informed by them cannot execute at that same close. The
target and the simulator share an execution lag of one day: decide after
the close of $t$, execute at the close of $t{+}1$, over strictly
consecutive trading days on which the name actually trades. This closes
the most common look-ahead leak; \texttt{execution\_lag}${=}0$ (same-close
MOC) is reserved for signal-decay diagnostics only.

\paragraph{Costs.} IDX round-trip frictions are modeled as $15$ bps
commission on the buy side and $25$ bps on the sell side (commission plus
sell tax), against a risk-free rate of $5.5\%$ (Sharpe is computed in
excess of it). The realistic simulator additionally charges each name's
own closing half-spread (default $35$ bps when the book is empty, capped
at $200$ bps); the training loss uses the effective commission-plus-typical-spread
number so it is not double-counted.

\section{Evaluation}
Evaluation is out-of-sample (walk-forward), never training loss. The
primary ranking metric is the annualized Sharpe ratio; supporting
metrics are cumulative return versus buy-and-hold IHSG (because
long-only carries beta), maximum drawdown, and turnover (a proxy for the
high real trading costs on IDX).

For robustness, each configuration is run over \textbf{eight random
seeds}, reporting mean $\pm$ standard deviation. Eight seeds provide a
stable variance estimate and support significance testing between
configurations; a Sharpe gap below roughly $0.1$ is likely seed noise
rather than a feature effect. The full study therefore comprises
$5 \times 8 = 40$ runs.

\section{Implementation Changes}
The study requires: (i) making \texttt{src/preprocess.py} config-driven,
replacing the fixed \texttt{FEATURE\_COLUMNS} constant with group
definitions and computing the feature superset while selecting active
groups from config; (ii) fixing the \texttt{open=0} bug for open-based
features; (iii) one YAML per experiment in \texttt{configs/} differing
only in \texttt{feature\_groups}; (iv) extending \texttt{compare.py} to
aggregate results into a single table (feature set $\to$ Sharpe $\to$
return $\to$ drawdown $\to$ turnover, mean $\pm$ std across seeds); and
(v) saving metrics and equity-curve plots to \texttt{results/}.

\section{Risks}
The main risks are overfitting from limited data combined with many
features (precisely the effect under test), feature collinearity (e.g.\
log\_volume vs.\ log\_value should not be read as new information),
seed noise (hence multi-seed runs), and look-ahead leakage in momentum,
rolling, and normalization steps (the most bug-prone area, to be covered
by anti-look-ahead unit tests). The long-only benchmark must remain
IHSG, not an absolute return figure.

\section{Definition of Done}
The study is complete when all five experiments (A--E) have run over
eight seeds each ($40$ runs), the comparison table is populated with
mean $\pm$ std for all metrics, the optimal feature set is identified
with justification (the Sharpe saturation point), and findings are
documented in \texttt{results/} and summarized in the project plan.

\end{document}
